% Replace appendix/experimental_cost.tex. Existing labels are retained.
\section{Computational cost and timing boundaries}
\label{app:parameter_counts}

We separate three questions: how many trainable parameters are used
per routing decision, how long the implemented local prediction paths
take, and how a warm specialist compares with a historical CFD run over
the same physical interval. Parameter accounting covers Circular H/P;
both timing studies cover Circular Periodic only. No Square or Pinball
runtime or complete hierarchical-pipeline speedup is inferred.

\subsection{Circular H/P: parameter utilization}

Total parameters count unique trainable parameters in each analyzed
specialist. Active parameters count unique trainable parameters used
by one batch-size-one prediction-path routing decision, not their union
across a batch, trajectory, or all Runge--Kutta stages. Included modules
are the encoder and refinement blocks, executed routers, the selected
group's shared experts, channel-specific Top-2 velocity and pressure
experts, and the learned pressure correction. Fixed POD bases, projected
operators, normalization statistics, and non-trainable buffers are excluded.

\begin{table}[tbp]
\centering
\caption{Circular trainable and per-decision active parameter counts.
Counts refer to unique trainable parameters, not FLOPs or memory usage.}
\label{tab:parameter_counts}
\begin{tabular}{@{}lrrr@{}}
\toprule
Specialist & Total trainable & Active prediction & Active fraction\\
\midrule
Hopf & 31,810,392 & 14,165,816 & 44.53\%\\
Periodic & 48,617,547 & 8,308,959 & 17.09\%\\
\bottomrule
\end{tabular}
\end{table}

Table~\ref{tab:parameter_counts} quantifies sparse use of the available
trainable capacity. It does not imply proportional reductions in storage,
operation count, or latency. The Dense control matches the active neural
parameter budget, not total model capacity or computational cost; it is
not identified with the older Vanilla-FNN-MoE, DataOnly-MoE, or Global MoE
configurations.

\subsection{Circular Periodic: local prediction latency}
\label{app:cost_local_latency}

\paragraph{Protocol and timing boundary.}
The native local-model comparison uses Circular Periodic seed 1248,
batch size one, $K=48$, three warm-up runs, and ten timed repetitions.
It includes modal rollout, feature generation, data transfers,
algebraic pressure reconstruction, all-step physical decoding, and
pressure gauge correction. It excludes physical-history encoding, E2
selection, descriptor construction, T2-C fusion, model/data loading,
and error computation. Thus it measures a local prediction path,
not a single neural forward pass or the complete hierarchy.

\begin{table}[tbp]
\centering\small
\caption{Native Circular Periodic local prediction cost at $K=48$.
Times are medians over ten repetitions after three warm-ups.
GPU memory denotes peak allocated bytes, not total device or host memory.}
\label{tab:local_latency}
\begin{tabular}{@{}lrrr@{}}
\toprule
Model & Median time (s) & GPU peak (bytes) & Time / Dense\\
\midrule
Proposed & 2.1848 & 204,238,336 & 3.32\\
Dense & 0.6579 & 43,111,936 & 1.00\\
Structured & 0.8652 & 43,316,736 & 1.32\\
\bottomrule
\end{tabular}
\end{table}

\paragraph{Interpretation.}
The proposed implementation is slower than both controls and has higher
peak allocated GPU memory in this measurement
(Table~\ref{tab:local_latency}). The native Galerkin NumPy path executes
on the CPU, but these aggregate timings do not isolate the contribution
of individual modules or establish the cause of the latency difference.
Sparse parameter activation therefore provides no measured runtime
advantage over Dense under this implementation and timing boundary.

\subsection{Circular Periodic: physical-input timing and historical FOM reference}
\label{app:cost_fom_reference}

\paragraph{Matched physical interval.}
A separate measurement of the original Periodic specialist (seed 1248)
adds full-size physical-history encoding and history right-hand-side
construction. At $\mathrm{Re}=70.314635$, the frozen first legal
$K=48$ window spans $t=1384.49597$ to $1983.19690$, a duration of
$598.70093$ physical time units. Its input consists of three
POD-reconstructed history frames on 97,368 points. All 48 predicted
velocity and pressure fields are decoded, with pressure gauge correction.
The model and input fields are resident before timing; acquiring the
history is not timed. This boundary differs from the local-model
comparison in Appendix~\ref{app:cost_local_latency}.

\paragraph{ROM measurement.}
On an RTX 4090 with a Xeon Platinum 8358P host, using one CPU numerical
thread, batch size one, three warm-ups, ten repetitions, and CUDA
synchronization, the median is $2.18747$ s. The mean is $2.18729$ s,
the sample standard deviation is $0.00462$ s, and the range is
$2.18034$--$2.19550$ s. Loading the model and inputs, writing predictions
to disk, E2 selection, descriptor construction, and T2-C fusion are
outside this measurement. Offline model construction and training are
also outside the online prediction timing.

\paragraph{Historical FOM measurement.}
The corresponding production OpenFOAM log contains 4,692 solver steps
and a ClockTime difference of 1,996 s. The largest discrepancy between
the ROM and FOM interval endpoints is $2.3\times10^{-5}$ physical time
units, attributable to stored timestamp precision. The production
solver used one process on a Core i5-14600K VMware host with 16 vCPUs
and at most five concurrent cases. This hardware information comes
from a subsequent source audit, not a contemporaneous controlled
benchmark. There is one FOM timing realization. The recorded interval
includes native solver field writes but excludes meshing, initialization,
periodic detection, and subsequent VTK/NPZ conversion.

\paragraph{Scope of the wall-time ratio.}
The ratio $1996/2.18747\simeq912.5$ compares a historical production
FOM interval with a warm, physical-input/physical-output local specialist.
It provides a case-specific online cost reference, not a controlled
end-to-end speedup: hardware, concurrency, repetition counts, and I/O
boundaries differ. In particular, the ROM measurement excludes outer
selection and fusion and cannot represent the cost of advancing and
assembling two specialists. A fully paired Top-2 end-to-end timing
under matched compute and I/O conditions remains unmeasured.
